Por tanto:
\[
\varepsilon=\frac{\sigma}{E}.
\]

Para una barra prismática:
\[
\sigma=\frac{N}{A}
\qquad\text{y}\qquad
\varepsilon=\frac{N}{EA}.
\]

Como
\[
\varepsilon=\frac{\dd u}{\dd x},
\]
obtenemos:
\[
\frac{\dd u}{\dd x}=\frac{N}{EA}.
\]

Si $N$, $E$ y $A$ son constantes:
\[
\Delta L=\frac{NL}{EA}.
\]

\safesection{Diagrama tensión--deformación}
El comportamiento típico de un acero dúctil puede representarse esquemáticamente como:

\begin{figure}[htbp]
\centering
\begin{tikzpicture}[x=0.95cm,y=0.82cm,>=Latex]
  % Ejes
  \draw[-{Latex[length=2.5mm]},BookGray,line width=.9pt] (0,0)--(10.4,0)
    node[below right,font=\sffamily\small] {$\varepsilon$};
  \draw[-{Latex[length=2.5mm]},BookGray,line width=.9pt] (0,0)--(0,6.35)
    node[above left,font=\sffamily\small] {$\sigma$};

  % Curva nominal de acero dúctil (esquema)
  \draw[BookBlue,line width=2.0pt]
    (0,0)
    -- (2.10,3.55)
    .. controls (2.35,3.82) and (2.65,3.62) .. (3.15,3.62)
    .. controls (4.20,3.62) and (5.55,4.60) .. (7.10,5.10)
    .. controls (7.75,5.28) and (8.60,4.60) .. (9.25,3.85);

  % Puntos y líneas auxiliares
  \fill[BookBlue] (2.10,3.55) circle (1.7pt);
  \fill[BookBlue] (7.10,5.10) circle (1.7pt);
  \fill[BookRed]  (9.25,3.85) circle (1.7pt);

  \draw[densely dashed,BookGray!70] (2.10,0)--(2.10,3.55);
  \draw[densely dashed,BookGray!70] (7.10,0)--(7.10,5.10);
  \draw[densely dashed,BookGray!70] (0,3.55)--(2.10,3.55);
  \draw[densely dashed,BookGray!70] (0,5.10)--(7.10,5.10);

  \node[left,font=\small] at (0,3.55) {$\sigma_y$};
  \node[left,font=\small] at (0,5.10) {$\sigma_u$};

  % Regiones
  \node[font=\sffamily\scriptsize,text=BookGray,align=center] at (1.15,-0.58)
    {zona\\elástica};
  \node[font=\sffamily\scriptsize,text=BookGray,align=center] at (2.80,-0.58)
    {fluencia};
  \node[font=\sffamily\scriptsize,text=BookGray,align=center] at (5.25,-0.58)
    {endurecimiento\\plástico};
  \node[font=\sffamily\scriptsize,text=BookGray,align=center] at (8.35,-0.58)
    {estricción};

  % Etiquetas de puntos
  \node[anchor=west,font=\sffamily\scriptsize,text=BookBlue] at (2.35,3.95)
    {inicio de fluencia};
  \node[anchor=west,font=\sffamily\scriptsize,text=BookBlue] at (7.25,5.38)
    {tensión última};
  \node[anchor=west,font=\sffamily\scriptsize,text=BookRed] at (9.35,3.82)
    {rotura};

  % Pendiente inicial
  \draw[BookGreen!75!black,-{Latex[length=1.8mm]},line width=.7pt]
    (1.05,2.45)--(1.68,3.48);
  \node[font=\sffamily\scriptsize,text=BookGreen!65!black] at (1.10,2.85)
    {$E=\mathrm{d}\sigma/\mathrm{d}\varepsilon$};
\end{tikzpicture}
\caption{Diagrama tensión-deformación nominal esquemático de un acero dúctil. La pendiente inicial representa el módulo de Young.}
\label{fig:diagrama-tension-deformacion}
\end{figure}

Podemos distinguir:
\begin{description}
  \item[Límite de proporcionalidad:] termina la proporcionalidad lineal.
  \item[Límite elástico:] por debajo de él no quedan deformaciones permanentes significativas.
  \item[Fluencia:] aparece deformación plástica apreciable.
  \item[Tensión última:] máximo de tensión nominal.
  \item[Estricción:] reducción localizada de la sección.
  \item[Rotura:] separación final de la probeta.
\end{description}

\begin{ojo}
Después de iniciarse la estricción, la \textbf{tensión nominal} puede disminuir porque sigue utilizando $A_0$. Eso no implica necesariamente que la tensión verdadera local también disminuya.
\end{ojo}

% ============================================================
\safechapter{El problema general de la elasticidad}
% ============================================================

La estructura lógica del problema elástico puede resumirse mediante tres grupos de ecuaciones.

\safesection{Equilibrio}
En el volumen:
\[
\nabla\cdot\mat{\sigma}+\vect b=\vect 0,
\]
donde $\vect b$ representa fuerzas por unidad de volumen.

En una parte de la frontera:
\[
\mat{\sigma}\vect n=\bar{\vect t}.
\]

\safesection{Compatibilidad / cinemática}
